\documentclass[11pt]{article}

\usepackage[T1]{fontenc}
\usepackage[utf8]{inputenc}
\usepackage{lmodern}
\usepackage[margin=1in]{geometry}
\usepackage{amsmath,amssymb,amsthm,mathtools}
\usepackage{microtype}
\usepackage{enumitem}
\usepackage{booktabs}
\usepackage{xcolor}
\usepackage{graphicx}
\usepackage[hidelinks]{hyperref}

\setlist{nosep,leftmargin=2em}
\allowdisplaybreaks

\newtheorem{theorem}{Theorem}[section]
\newtheorem{proposition}[theorem]{Proposition}
\newtheorem{lemma}[theorem]{Lemma}
\newtheorem{corollary}[theorem]{Corollary}
\theoremstyle{definition}
\newtheorem{definition}[theorem]{Definition}
\newtheorem{remark}[theorem]{Remark}
\newtheorem{question}[theorem]{Question}

\newcommand{\N}{\lVert\,\cdot\,\rVert}
\newcommand{\norm}[1]{\left\lVert #1\right\rVert}
\newcommand{\qnorm}[1]{\left\lvert #1\right\rvert_q}
\newcommand{\rnorm}[1]{\left\lvert #1\right\rvert_r}
\newcommand{\ellTwo}{\ell_2}
\newcommand{\Hilb}{\mathsf H}
\newcommand{\dist}{d_{\mathrm{BM}}}
\newcommand{\ess}{\mathrm{ess}}
\newcommand{\spec}{\sigma}
\newcommand{\Id}{I}
\newcommand{\Kcal}{\mathcal K}
\newcommand{\Bcal}{\mathcal B}
\newcommand{\Calk}{\mathcal C}
\newcommand{\R}{\mathbb R}
\newcommand{\C}{\mathbb C}
\newcommand{\F}{\mathbb F}

\title{\textbf{The Isometrically Homogeneous Hilbert-Space Problem}\\[0.4em]
\large Audit of the Supplied Argument, Literature Status, and Structural Reductions}
\author{Research report prepared from the supplied source paper and partial solution}
\date{22 June 2026}

\begin{document}
\maketitle

\begin{center}
\fbox{\begin{minipage}{0.91\textwidth}
\textbf{Status.} The supplied conditional theorem is correct.  Its stated
``symmetric distortion'' corollary needs a small square-root normalization in
its proof, but its qualitative conclusion is correct.  I found no published
full proof or counterexample in the literature search described below.  The
problem therefore remains open on the evidence available to this report.
Sections~\ref{sec:optimal}--\ref{sec:shift} give additional rigorous reductions:
an optimal Hilbert position exists; every infinite-dimensional subspace attains
the full optimal distortion range; optimal Hilbert structures are unique modulo
compact operators; every linear isometric embedding is a compact perturbation
of a Hilbert isometry; and a hypothetical counterexample admits a pure
unilateral-shift normal form.
\end{minipage}}
\end{center}

\begin{abstract}
Let $X$ be a separable infinite-dimensional Banach space linearly isometric to
every one of its closed infinite-dimensional subspaces.  C\'uth, Dole\v{z}al,
Doucha and Kurka asked whether $X$ must be isometric to $\ell_2$.  The
isomorphic homogeneous-space theorem implies only that $X$ is isomorphic to
$\ell_2$, and arbitrary distortion of Hilbert space prevents the usual shortcut
from isomorphism to almost-isometry.

This report first verifies the supplied almost-Hilbertian stabilization
argument.  It then records the literature status and develops a sequence of
intrinsic reductions.  If $D=d_{\mathrm{BM}}(X,\ell_2)$, the infimum defining
$D$ is attained by an equivalent Hilbert norm $|\cdot|_q$ normalized by
$|x|_q\leq\|x\|\leq D|x|_q$.  In a non-Hilbert counterexample, every closed
infinite-dimensional subspace has norm-ratio infimum $1$ and supremum $D$ in
this same Hilbert position.  Any other normalized optimal Hilbert norm differs
from $q$ by a compact positive perturbation.  Consequently, every linear
$\|\cdot\|$-isometric embedding $U:X\to X$ satisfies $U^*U-I$ compact with
respect to $q$, and hence is a compact perturbation of a Hilbert isometric
embedding.  Finally, after passing to an isometric subspace, one may arrange
that $X$ is an equivalent renorming of $\ell_2$ for which the canonical right
shift is simultaneously a Banach-space isometry and a Hilbert isometry.
These conclusions sharply constrain a counterexample but do not yet force
$D=1$.
\end{abstract}

\tableofcontents

\section*{Packet status and source evidence}
\addcontentsline{toc}{section}{Packet status and source evidence}

This is a partial-result packet for arXiv:2204.06834, Question 3 of
C\'uth--Dole\v{z}al--Doucha--Kurka.  It does not claim a full solution or a
counterexample.  The source problem remains open on the evidence reviewed here.
The contribution of this upgraded packet is to preserve the verified
almost-Hilbertian obstruction from the previous packet, correct the
square-root normalization in its symmetric-distortion corollary, and record
additional reductions that any non-Hilbert counterexample must satisfy.

\begin{itemize}
\item Previous active packet preserved at
\path{history/2026-06-22_previous_almost_hilbertian_partial/}.
\item Supplied upgrade report preserved in the dated evidence folder.
\item Source paper preserved as \path{source_paper.pdf}.
\end{itemize}

\clearpage
\begin{center}
\includegraphics[width=\textwidth]{figures/open_problem_crop.png}
\end{center}
\clearpage

\section{The problem and terminology}

Throughout, $\F$ denotes either $\R$ or $\C$.  A Banach space
$X=(X,\|\cdot\|)$ is called \emph{isometrically homogeneous} in this report if,
for every closed infinite-dimensional subspace $Y\subseteq X$, there is a
surjective linear isometry $U:X\to Y$.

The question considered by C\'uth, Dole\v{z}al, Doucha and Kurka is the
following.

\begin{question}[Isometric homogeneous-space problem]\label{q:main}
If a separable infinite-dimensional Banach space is isometrically homogeneous,
must it be linearly isometric to $\ell_2$?
\end{question}

The homogeneous-space theorem of Komorowski--Tomczak-Jaegermann and Gowers
implies that such an $X$ is linearly \emph{isomorphic} to $\ell_2$.  Thus one
may fix an equivalent Hilbert norm on the same vector space.  The open issue is
whether the given norm itself satisfies the parallelogram identity.

For two isomorphic Banach spaces $E,F$, write
\[
 d_{\mathrm{BM}}(E,F)
 =\inf\{\|T\|\,\|T^{-1}\|:T:E\to F\text{ is a linear isomorphism}\}.
\]
When $F$ is a separable infinite-dimensional Hilbert space, we abbreviate
$d_{\mathrm{BM}}(E,F)$ to $d(E,\ell_2)$.

\section{Literature status}\label{sec:literature}

The journal version of the source paper appeared in 2024 and still lists
Question~\ref{q:main} as open \cite{CDDK}.  The strongest directly targeted
result I found is due to de Rancourt:

\begin{theorem}[de Rancourt \cite{deRancourt}]\label{thm:derancourt}
If $X$ is isometrically homogeneous, then for every $\varepsilon>0$ the space
$X$ admits a $(1+\varepsilon)$-unconditional basis.
\end{theorem}

De Rancourt explicitly identifies the missing ingredient as a sufficiently
strong quantitative form of the Komorowski--Tomczak-Jaegermann theorem.  The
obstruction is genuine.  Odell and Schlumprecht proved that $\ell_2$ is
arbitrarily distortable: for every $\lambda>1$ there is an equivalent norm on
$\ell_2$ whose restriction to every infinite-dimensional subspace has
oscillation exceeding $\lambda$ relative to the original Hilbert norm
\cite{OdellSchlumprecht}.  They also constructed, for prescribed $\lambda>1$,
equivalent norms with no $\lambda$-unconditional basic sequence.  Thus
Theorem~\ref{thm:derancourt} rules out the most hostile known distortions, but
it does not by itself supply an almost-Euclidean infinite-dimensional
restriction.

A second warning comes from Casazza's survey \cite{Casazza}: exact recurrence
of all finite-dimensional isometry types in every infinite-dimensional
subspace is not sufficient in arbitrary Banach spaces; $c_0$ has that
recurrence property.  Any proof of Question~\ref{q:main} must therefore use the
additional Hilbertian isomorphism structure, not only finite-dimensional
recurrence.

The search performed for this report covered the exact wording of the
question, variants of ``isometrically homogeneous Banach space'', the title
and citations of the source paper, citations of de Rancourt's paper, and
arXiv/journal material through 22 June 2026.  No later full proof or
counterexample was located.  This is evidence about the indexed literature,
not a claim that an unpublished or obscure result cannot exist.

\section{Audit of the supplied partial result}\label{sec:audit}

The central theorem in the supplied packet is correct.  It is useful to state
it in a slightly normalized form.

\begin{theorem}[Verified supplied theorem]\label{thm:supplied}
Let $X=(X,\|\cdot\|)$ be separable, infinite-dimensional and isometrically
homogeneous.  Suppose that there is an equivalent Hilbert norm $|\cdot|$ such
that for every $\varepsilon>0$ there are a closed infinite-dimensional
subspace $Y\subseteq X$ and $a>0$ satisfying
\[
 a|y|\leq \|y\|\leq (1+\varepsilon)a|y|
 \qquad(y\in Y).
\]
Then $X$ is linearly isometric to $\ell_2$.
\end{theorem}

\begin{proof}
Fix a finite-dimensional subspace $F\subseteq X$ and $\varepsilon>0$.  Choose
$Y$ and $a$ as in the hypothesis, and let $U:X\to Y$ be a surjective linear
$\|\cdot\|$-isometry.  For $x\in F$,
\[
 a|Ux|\leq \|Ux\|=\|x\|\leq(1+\varepsilon)a|Ux|.
\]
Since $(U(F),|\cdot|)$ is Euclidean,
$d_{\mathrm{BM}}(F,\ell_2^{\dim F})\leq1+\varepsilon$.  Letting
$\varepsilon\downarrow0$ gives distance one.  In finite dimension, distance
one is attained: after normalizing almost-isometries, compactness of the unit
sphere in the relevant operator space yields an actual linear isometry.
Therefore every finite-dimensional subspace of $X$ is Euclidean.  In
particular every two-dimensional subspace satisfies the parallelogram law, so
the Jordan--von Neumann theorem makes $\|\cdot\|$ an inner-product norm.
Separable infinite-dimensional Hilbert spaces are linearly isometric to
$\ell_2$.
\end{proof}

There is one normalization point in the packet's corollary.  For a fixed
Hilbert norm $|\cdot|$ and an infinite-dimensional subspace $Y$, put
\[
 m_Y=\inf_{|y|=1,\,y\in Y}\|y\|,
 \qquad
 M_Y=\sup_{|y|=1,\,y\in Y}\|y\|.
\]
Then
\begin{equation}\label{eq:symmetric-distortion}
 \inf_{a>0}\sup_{0\ne y\in Y}
 \max\left\{\frac{\|y\|}{a|y|},\frac{a|y|}{\|y\|}\right\}
 =\sqrt{\frac{M_Y}{m_Y}}.
\end{equation}
Indeed the left side is $\inf_{a>0}\max\{M_Y/a,a/m_Y\}$, minimized at
$a=\sqrt{M_Ym_Y}$.  The theorem's hypothesis is the ratio condition
$M_Y/m_Y\leq1+\varepsilon$, whereas the symmetric expression in
\eqref{eq:symmetric-distortion} is its square root.

\begin{corollary}[Corrected proof of the packet's corollary]\label{cor:packet}
If a non-Hilbert counterexample to Question~\ref{q:main} exists, then for every
equivalent Hilbert norm $|\cdot|$ there is $\delta>0$ such that every closed
infinite-dimensional $Y\subseteq X$ satisfies
\[
 \inf_{a>0}\sup_{0\ne y\in Y}
 \max\left\{\frac{\|y\|}{a|y|},\frac{a|y|}{\|y\|}\right\}
 \geq1+\delta.
\]
\end{corollary}

\begin{proof}
Negating the stabilization hypothesis in Theorem~\ref{thm:supplied} gives an
$\eta>0$ such that $M_Y/m_Y>1+\eta$ for every infinite-dimensional $Y$.
Equation~\eqref{eq:symmetric-distortion} then gives a uniform lower bound
strictly larger than $\sqrt{1+\eta}$; decreasing the bound slightly yields the
stated non-strict inequality.  Thus the qualitative conclusion in the packet
is correct, but the square-root conversion should be made explicit.
\end{proof}

\section{An intrinsic almost-Hilbert criterion}\label{sec:intrinsic}

The fixed ambient Hilbert norm in Theorem~\ref{thm:supplied} is unnecessary.
The real mechanism is intrinsic Banach--Mazur stabilization.

\begin{theorem}[Intrinsic criterion]\label{thm:intrinsic}
Let $X$ be isometrically homogeneous.  If
\[
 \inf\{d(Y,\ell_2):Y\subseteq X\text{ closed and infinite-dimensional}\}=1,
\]
then $X$ is linearly isometric to $\ell_2$.
\end{theorem}

\begin{proof}
Fix a finite-dimensional $F\subseteq X$ and $\varepsilon>0$.  Choose an
infinite-dimensional $Y\subseteq X$ with $d(Y,\ell_2)<1+\varepsilon$, and let
$U:X\to Y$ be a surjective $\|\cdot\|$-isometry.  Restricting an
$(1+\varepsilon)$-Hilbertian structure on $Y$ to $U(F)$ shows
$d(F,\ell_2^{\dim F})<1+\varepsilon$.  The rest is exactly the compactness and
parallelogram argument from Theorem~\ref{thm:supplied}.
\end{proof}

\begin{remark}
Theorem~\ref{thm:intrinsic} isolates the precise gap left by arbitrary
distortion.  A full positive solution would follow from proving that an
isometrically homogeneous renorming of $\ell_2$ cannot remain a fixed positive
distance from Hilbert space on every infinite-dimensional subspace.
\end{remark}

\section{Optimal Hilbert position}\label{sec:optimal}

From now on let $X=(H,\|\cdot\|)$ be isomorphic to a separable
infinite-dimensional Hilbert space, and set
\[
 D=d(X,\ell_2).
\]
The first useful point is that, although Banach--Mazur infima need not be
attained for arbitrary pairs of infinite-dimensional spaces, the Hilbert
target permits attainment.

\begin{theorem}[Attainment of the Hilbert distance]\label{thm:attainment}
There is an equivalent Hilbert norm $|\cdot|_q$ on $H$ such that
\begin{equation}\label{eq:optimal-position}
 |x|_q\leq\|x\|\leq D|x|_q\qquad(x\in H).
\end{equation}
In particular, if $D=1$, then $\|\cdot\|$ is a Hilbert norm.
\end{theorem}

\begin{proof}
Fix a reference Hilbert norm $|\cdot|_0$.  For each $k$ choose a Hilbert norm
$|\cdot|_k$, scaled so that
\[
 |x|_k\leq\|x\|\leq(D+k^{-1})|x|_k.
\]
Write $|x|_k^2=\langle A_kx,x\rangle_0$ with $A_k$ positive and invertible.
Equivalence of $|\cdot|_0$ and $\|\cdot\|$, together with the displayed
inequalities, gives constants $0<c<C<\infty$ independent of $k$ such that
$cI\leq A_k\leq CI$.  The order interval $[cI,CI]$ is compact in the weak
operator topology after passing to a subnet (or, on separable $H$, a
subsequence).  Let $A_k\to A$ weakly.  Then $cI\leq A\leq CI$ and
$|x|_q^2:=\langle Ax,x\rangle_0$ defines an equivalent Hilbert norm.  Passing
to the limit pointwise in the two inequalities gives
\eqref{eq:optimal-position}.
\end{proof}

Now impose isometric homogeneity.

\begin{theorem}[Hereditary full distortion]\label{thm:endpoints}
Assume $X$ is isometrically homogeneous and let $q$ be normalized as in
\eqref{eq:optimal-position}.  For every closed infinite-dimensional
$Y\subseteq H$,
\begin{equation}\label{eq:endpoints}
 \inf_{\substack{y\in Y\\|y|_q=1}}\|y\|=1,
 \qquad
 \sup_{\substack{y\in Y\\|y|_q=1}}\|y\|=D.
\end{equation}
\end{theorem}

\begin{proof}
Let the two quantities be $m_Y$ and $M_Y$.  From
\eqref{eq:optimal-position}, $1\leq m_Y\leq M_Y\leq D$.  Since $Y$ is
isometric to $X$, $d(Y,\ell_2)=D$.  The identity between $(Y,|\cdot|_q)$ and
$(Y,\|\cdot\|)$ has distortion $M_Y/m_Y$, so
\[
 D\leq M_Y/m_Y\leq D.
\]
Hence $M_Y=Dm_Y$.  The bounds $M_Y\leq D$ and $m_Y\geq1$ force
$m_Y=1$ and $M_Y=D$.
\end{proof}

\begin{corollary}\label{cor:counter-distortion}
A non-Hilbert counterexample has $D>1$ and, in an optimal Hilbert position,
its entire distortion interval $[1,D]$ is unavoidable on every
infinite-dimensional subspace.  In particular, neither endpoint can be
removed by passing to a subspace.
\end{corollary}

This is stronger than the supplied corollary because it uses an optimal
Hilbert norm and identifies the exact hereditary endpoint values.

\section{Rigidity of Hilbert structures modulo compact operators}
\label{sec:compact}

Let $q$ be an optimal Hilbert norm as above.  If $r$ is another Hilbert norm,
write
\[
 |x|_r^2=\langle Ax,x\rangle_q
\]
for its positive invertible comparison operator $A$.  The next result converts
hereditary full distortion into an essential-spectrum constraint.

\begin{theorem}[Near-optimal essential-spectrum rigidity]\label{thm:ess-spectrum}
Assume $X$ is isometrically homogeneous and $q$ satisfies
\eqref{eq:optimal-position}.  Let $r$ be a Hilbert norm normalized by
\[
 |x|_r\leq\|x\|\leq D'|x|_r
 \qquad(x\in H),
\]
where $D'\geq D$.  If $|x|_r^2=\langle Ax,x\rangle_q$, then
\begin{equation}\label{eq:ess-spectrum}
 \sigma_{\mathrm{ess}}(A)\subseteq
 \left[\left(\frac{D}{D'}\right)^2,1\right].
\end{equation}
\end{theorem}

\begin{proof}
Suppose first that the essential spectrum meets $(1,\infty)$.  Then for some
$\eta>0$ the spectral subspace
$Y=\mathbf 1_{[1+\eta,\infty)}(A)H$ is infinite-dimensional.  On $Y$,
$|y|_r\geq\sqrt{1+\eta}|y|_q$, and $|y|_r\leq\|y\|$.  Therefore
$\|y\|/|y|_q\geq\sqrt{1+\eta}$ on $Y\setminus\{0\}$, contradicting the first
identity in \eqref{eq:endpoints}.

Now suppose the essential spectrum meets
$\bigl(0,(D/D')^2\bigr)$.  Choose $c<D/D'$ so that
$Y=\mathbf 1_{(0,c^2]}(A)H$ is infinite-dimensional.  On $Y$,
$|y|_r\leq c|y|_q$, whence
\[
 \|y\|\leq D'|y|_r\leq D'c|y|_q<D|y|_q,
\]
contradicting the second identity in \eqref{eq:endpoints}.  This proves
\eqref{eq:ess-spectrum}.
\end{proof}

\begin{corollary}[Compact uniqueness of optimal Hilbert positions]
\label{cor:compact-unique}
If $q_1$ and $q_2$ are both normalized optimal Hilbert norms,
\[
 |x|_{q_i}\leq\|x\|\leq D|x|_{q_i}\quad(i=1,2),
\]
and $|x|_{q_1}^2=\langle Ax,x\rangle_{q_2}$, then $A-I$ is compact.
\end{corollary}

\begin{proof}
Apply Theorem~\ref{thm:ess-spectrum} with $D'=D$.  It gives
$\sigma_{\mathrm{ess}}(A)\subseteq\{1\}$.  The essential spectrum of a
bounded self-adjoint operator on an infinite-dimensional Hilbert space is
nonempty, so it equals $\{1\}$.  Equivalently, $A-I$ is compact.
\end{proof}

\section{Isometries are Hilbertian modulo compacts}\label{sec:isometries}

The preceding compact uniqueness applies directly to every into-isometry.
Adjoints and compactness in this section are taken with respect to the optimal
Hilbert norm $q$.

\begin{theorem}[Compact polar rigidity]\label{thm:isom-compact}
Let $U:X\to X$ be a linear $\|\cdot\|$-isometric embedding, not necessarily
surjective.  Then
\begin{equation}\label{eq:ustaru}
 U^*U-I\quad\text{is compact}.
\end{equation}
Consequently, in the polar decomposition $U=V|U|$, the operator $V$ is a
$q$-Hilbert isometric embedding with the same range as $U$, and
\[
 U-V\quad\text{is compact}.
\]
If $U$ is surjective, $V$ is unitary.
\end{theorem}

\begin{proof}
Define $|x|_{q_U}=|Ux|_q$.  Since $U$ is a $\|\cdot\|$-isometry,
\[
 |x|_{q_U}=|Ux|_q\leq\|Ux\|=\|x\|
 \leq D|Ux|_q=D|x|_{q_U}.
\]
Thus $q_U$ is another normalized optimal Hilbert norm.  Its comparison
operator with $q$ is $U^*U$, so Corollary~\ref{cor:compact-unique} gives
\eqref{eq:ustaru}.  Functional calculus yields
$|U|=(U^*U)^{1/2}=I+K$ with $K$ compact.  Hence
$U=V(I+K)$ and $U-V=VK$ is compact.  Since $|U|$ is invertible, $U$ and $V$
have the same range.
\end{proof}

\begin{proposition}[Dual quotient homogeneity]\label{prop:dual-quotient}
If $X$ is isometrically homogeneous, then every infinite-dimensional quotient
of $X^*$ is linearly isometric to $X^*$.
\end{proposition}

\begin{proof}
The homogeneous-space theorem gives $X\simeq\ell_2$, so $X$ is reflexive.
Let $Z\subseteq X^*$ be closed and suppose $X^*/Z$ is infinite-dimensional.
Because $X$ is reflexive, norm-closed subspaces of $X^*$ are weak-star closed.
With $Y=Z_\perp\subseteq X$, one has $Z=Y^\perp$, and restriction gives the
canonical quotient isometry
\[
 X^*/Z=X^*/Y^\perp\cong Y^*.
\]
The quotient is infinite-dimensional exactly when $Y$ is.  Homogeneity gives
$Y\cong X$ isometrically, hence $Y^*\cong X^*$ isometrically.
\end{proof}

\section{Interaction with near-unconditional bases}\label{sec:unconditional}

Theorem~\ref{thm:derancourt} can be combined with
Theorem~\ref{thm:ess-spectrum} to obtain a precise ``unitary modulo compacts''
form of the sign symmetries.  This does not solve the problem because the basis
may vary with $\varepsilon$, but it narrows the remaining quantitative gap.

\begin{proposition}[Essential Hilbertianity of sign changes]
\label{prop:signs}
Let $(x_n)$ be a $K$-unconditional basis of $X$, using the sign-change
convention: every coordinate sign operator $S_\theta$ satisfies
$\|S_\theta\|_{X\to X}\leq K$.  Relative to an optimal Hilbert norm $q$,
\begin{equation}\label{eq:sign-spectrum}
 \sigma_{\mathrm{ess}}(S_\theta^*S_\theta)
 \subseteq[K^{-2},K^2].
\end{equation}
Moreover, there is a $q$-unitary $W_\theta$ such that
\begin{equation}\label{eq:essential-distance}
 \|S_\theta-W_\theta\|_{\mathrm{ess},q}\leq K-1.
\end{equation}
\end{proposition}

\begin{proof}
Since $S_\theta^{-1}=S_\theta$, unconditionality gives
$K^{-1}\|x\|\leq\|S_\theta x\|\leq K\|x\|$.  Define
\[
 |x|_r=K^{-1}|S_\theta x|_q.
\]
Then $|x|_r\leq\|x\|$ and
\[
 \|x\|\leq K\|S_\theta x\|
 \leq KD|S_\theta x|_q=K^2D|x|_r.
\]
The comparison operator of $r$ with $q$ is
$A=K^{-2}S_\theta^*S_\theta$.  Applying
Theorem~\ref{thm:ess-spectrum} with $D'=K^2D$ gives
$\sigma_{\mathrm{ess}}(A)\subseteq[K^{-4},1]$, which is
\eqref{eq:sign-spectrum}.

Let $S_\theta=W_\theta|S_\theta|$ be its Hilbert polar decomposition.  Because
$S_\theta$ is invertible, $W_\theta$ is unitary.  Equation
\eqref{eq:sign-spectrum} implies
$\sigma_{\mathrm{ess}}(|S_\theta|)\subseteq[K^{-1},K]$, so
\[
 \||S_\theta|-I\|_{\mathrm{ess},q}
 \leq\max\{K-1,1-K^{-1}\}=K-1.
\]
Finally $S_\theta-W_\theta=W_\theta(|S_\theta|-I)$.
\end{proof}

\begin{corollary}
For an isometrically homogeneous space and every $\varepsilon>0$, there is a
basis whose entire sign-change group is, in the Calkin algebra of the optimal
Hilbert space, within $\varepsilon$ of unitary operators.
\end{corollary}

\begin{remark}
The unresolved step is to turn these basis-dependent asymptotic unitary
symmetries into a single infinite-dimensional subspace on which
$\|\cdot\|/|\cdot|_q$ has oscillation tending to one.  Arbitrary distortion
shows that this conclusion cannot be obtained from the Hilbertian isomorphism
alone, while Proposition~\ref{prop:signs} shows exactly what additional
operator-theoretic information de Rancourt's theorem provides.
\end{remark}

\section{A pure unilateral-shift normal form}\label{sec:shift}

A counterexample, if it exists, may be put in a particularly rigid self-similar
position.  The construction uses a shift-invariant Banach limit and the Wold
decomposition.

\begin{theorem}[Shift normal form]\label{thm:shift}
Let $X$ be isometrically homogeneous.  There is a Banach space isometric to
$X$, represented as $(\ell_2,\|\cdot\|)$, and an optimal Hilbert norm equal to
the standard $\ell_2$ norm, such that the canonical right shift
\[
 S(a_1,a_2,\ldots)=(0,a_1,a_2,\ldots)
\]
is simultaneously a $\|\cdot\|$-isometry and a Hilbert isometry.  The
represented space remains isometrically homogeneous.
\end{theorem}

\begin{proof}
Choose a closed hyperplane $H_1\subset X$ and a surjective
$\|\cdot\|$-isometry $T:X\to H_1$, viewed as an into-isometry $T:X\to X$.
Start with an optimal Hilbert norm $|\cdot|_0$ satisfying
$|x|_0\leq\|x\|\leq D|x|_0$.  Since $T$ is a Banach-space isometry, for every
$n\geq0$,
\[
 D^{-1}|x|_0\leq |T^nx|_0\leq D|x|_0.
\]
Let $\operatorname{LIM}$ be a shift-invariant Banach limit and define
\[
 |x|_T^2=\operatorname{LIM}_{n\to\infty}|T^nx|_0^2.
\]
This is a positive definite quadratic form equivalent to $|\cdot|_0$; the
parallelogram identity passes through the Banach limit.  Shift invariance gives
$|Tx|_T=|x|_T$.  More precisely, because
$|T^nx|_0\leq\|T^nx\|=\|x\|$ and
$\|x\|=\|T^nx\|\leq D|T^nx|_0$, one has
\[
 |x|_T\leq\|x\|\leq D|x|_T.
\]
Thus $|\cdot|_T$ is still optimal.

Apply the Wold decomposition to the Hilbert isometry $T$:
\[
 H=H_u\oplus\bigoplus_{n\geq0}T^nW,
 \qquad W=H\ominus T(H).
\]
Since $T(H)$ has codimension one, $\dim W=1$.  The shift part
$H_s=\bigoplus_{n\geq0}T^nW$ is infinite-dimensional and closed for both
$|\cdot|_T$ and $\|\cdot\|$.  By homogeneity, $(H_s,\|\cdot\|)$ is isometric
to $X$.  Replacing $X$ by this isometric copy removes the unitary part.  If
$w$ is a Hilbert unit vector spanning $W$, the sequence $(T^nw)_{n\geq0}$ is an
orthonormal basis of $H_s$, and under the resulting identification with
$\ell_2$, $T$ becomes the canonical right shift.  Any infinite-dimensional
closed subspace of $H_s$ is also one of $X$, so homogeneity is preserved.
\end{proof}

\section{The resulting counterexample profile}\label{sec:profile}

Combining the preceding results gives a compact description of what remains to
be ruled out.

\begin{theorem}[Necessary profile of a counterexample]\label{thm:profile}
If Question~\ref{q:main} has a negative answer, then there is an equivalent norm
$\|\cdot\|$ on $\ell_2$ and a number $D>1$ with all of the following properties.
\begin{enumerate}[label=\textup{(\roman*)}]
\item $|x|_2\leq\|x\|\leq D|x|_2$, and $D=d((\ell_2,\|\cdot\|),\ell_2)$.
\item On every infinite-dimensional closed $Y\subseteq\ell_2$,
\[
 \inf_{|y|_2=1,\,y\in Y}\|y\|=1,
 \qquad
 \sup_{|y|_2=1,\,y\in Y}\|y\|=D.
\]
\item The canonical right shift is a $\|\cdot\|$-isometry.
\item Every linear $\|\cdot\|$-isometric embedding $U$ is a compact
perturbation of a Hilbert isometric embedding; equivalently,
$U^*U-I$ is compact.
\item For every $\varepsilon>0$, the space has a
$(1+\varepsilon)$-unconditional basis, and the associated sign changes satisfy
Proposition~\ref{prop:signs}.
\item Every finite-dimensional subspace of $X$ embeds linearly isometrically
into every infinite-dimensional subspace of $X$.
\item Every infinite-dimensional quotient of $X^*$ is linearly isometric to
$X^*$.
\end{enumerate}
\end{theorem}

\begin{proof}
Items (i)--(iv) are Theorems~\ref{thm:attainment},
\ref{thm:endpoints}, \ref{thm:isom-compact} and \ref{thm:shift}.  Item (v) is
de Rancourt's theorem together with Proposition~\ref{prop:signs}.  For (vi), if
$F\subset X$ is finite-dimensional and $Y\subset X$ is infinite-dimensional,
a surjective isometry $U:X\to Y$ maps $F$ isometrically into $Y$.  Item (vii)
is Proposition~\ref{prop:dual-quotient}.
\end{proof}

\section{Why the remaining step is difficult}\label{sec:bottlenecks}

The reductions above leave a sharply formulated target: prove that the profile
in Theorem~\ref{thm:profile} is impossible when $D>1$.  Several tempting
shortcuts fail.

\begin{enumerate}[label=\textbf{B\arabic*.}]
\item \textbf{Isomorphism does not give an almost-Hilbert subspace.}
Odell--Schlumprecht distortion directly blocks this inference.

\item \textbf{A $1$-unconditional basis does not itself make the norm
Euclidean.}  For example,
\[
 \|x\|_\alpha=\|x\|_2+\alpha\|x\|_\infty
 \qquad(\alpha>0)
\]
is an equivalent non-Hilbert norm on $\ell_2$ with a $1$-unconditional
canonical basis.  It is not a counterexample: suitable flat block subspaces
are arbitrarily close to Hilbert.  The missing issue is to derive such a flat
subspace uniformly from the hereditary self-isometry property.

\item \textbf{Finite-dimensional recurrence is insufficient by itself.}
Casazza's $c_0$ observation shows that every finite-dimensional isometry type
can recur in every infinite-dimensional subspace without the ambient norm
being Hilbertian.

\item \textbf{Compact perturbation is not yet exact rigidity.}
The Calkin algebra contains many proper isometries and many equivalent
infinite projections.  Thus Theorem~\ref{thm:isom-compact} is a substantial
operator reduction, but Fredholm or Calkin-algebra information alone does not
force $D=1$.

\item \textbf{The shift symmetry is one-sided.}  Shift invariance gives exact
self-similarity under insertion of a zero coordinate, but not invariance under
arbitrary permutations or Hilbert rotations.  A proof must manufacture enough
additional symmetry from homogeneity, or show that hereditary full distortion
is incompatible with the available near-unitary sign symmetries.
\end{enumerate}

A concrete next lemma that would finish the problem is the following.

\begin{question}[Quantitative stabilization target]\label{q:target}
Does every equivalent norm on $\ell_2$ satisfying items (ii), (iii), and (v) of
Theorem~\ref{thm:profile} have, for every $\varepsilon>0$, an
infinite-dimensional subspace $Y$ with
$d(Y,\ell_2)<1+\varepsilon$?
\end{question}

A positive answer to Question~\ref{q:target}, even without using all the other
items in the profile, would combine with Theorem~\ref{thm:intrinsic} to solve
the original problem.

\section{Conclusion}

The supplied theorem is sound and identifies a valid sufficient condition.
The only correction needed is to distinguish the ratio $M_Y/m_Y$ from the
symmetric distortion $\sqrt{M_Y/m_Y}$ in the corollary.

No full answer or counterexample was found in the literature search, and the
2024 journal source still presents the question as open.  The additional
results in this report do not settle it, but they reduce a hypothetical
counterexample to a highly constrained renorming of $\ell_2$: it must be
hereditarily fully distortive in an attained optimal Hilbert position, possess
near-perfect unconditional bases, have all isometric embeddings Hilbertian
modulo compact operators, and admit an exact pure unilateral-shift symmetry.
The remaining bottleneck is to turn these asymptotic and operator-theoretic
symmetries into an almost-Hilbertian infinite-dimensional restriction.

\medskip
\noindent\textbf{Novelty caution.}  The deductions in
Sections~\ref{sec:optimal}--\ref{sec:shift} were developed and checked for this
report.  I did not find these exact formulations in the searched literature,
but they have not undergone peer review.

\begin{thebibliography}{99}

\bibitem{Casazza}
P.~G. Casazza,
\emph{Some questions arising from the homogeneous Banach space problem},
arXiv:math/9205207 (1992).

\bibitem{CDDK}
M. C\'uth, M. Dole\v{z}al, M. Doucha, and O. Kurka,
\emph{Polish spaces of Banach spaces: complexity of isometry and isomorphism
classes},
J. Inst. Math. Jussieu \textbf{23} (2024), no.~4, 1919--1957;
\href{https://doi.org/10.1017/S1474748023000440}{doi:10.1017/S1474748023000440}.

\bibitem{deRancourt}
N. de Rancourt,
\emph{Spectral-free methods in the theory of hereditarily indecomposable
Banach spaces},
Bull. Belg. Math. Soc. Simon Stevin \textbf{27} (2020), 775--787;
\href{https://doi.org/10.36045/j.bbms.200123}{doi:10.36045/j.bbms.200123}.

\bibitem{Gowers}
W.~T. Gowers,
\emph{An infinite Ramsey theorem and some Banach-space dichotomies},
Ann. of Math. (2) \textbf{156} (2002), 797--833.

\bibitem{KTJ}
R.~A. Komorowski and N. Tomczak-Jaegermann,
\emph{Banach spaces without local unconditional structure},
Israel J. Math. \textbf{89} (1995), 205--226;
erratum, Israel J. Math. \textbf{105} (1998), 85--92.

\bibitem{OdellSchlumprecht}
E. Odell and Th. Schlumprecht,
\emph{The distortion problem},
Acta Math. \textbf{173} (1994), 259--281.

\end{thebibliography}

\end{document}
